\section{Affine and logarithmic calculations}
\label{app:calculations}

This appendix collects the basis conventions and local algebra used in the
main proofs.  It records the lattice data, the widths of the monodromies
and the separation estimate in explicit coordinates.  The conventions for
the two sides of the discrete Legendre transform are those of
Section~\ref{subsec:parity-cover}: frames and developments are in the plane
$M^\vee\otimes\RR$ of $\Delta$, and monodromies of characters act on the fan
lattice $M$.  Elements of $M$ and of $M^\vee$ are written in standard
coordinates, which on $M^\vee\otimes\RR$ are the coordinates $(s_1,s_2)$ of
\eqref{eq:polygons}; elements of $\widetilde M$ and $\widetilde M^\vee$ are
written in $F$-coordinates (Section~\ref{subsec:exponent-lattices}).

\subsection{Straight-boundary frames}

At a vertex $p_V$ of $\Delta$, let $\ell_V,r_V$ be the primitive boundary
tangents toward its two neighbors, and let $-p_V$ be the inward radial
vector.  Both $(\ell_V,-p_V)$ and $(r_V,-p_V)$ are unimodular.  Define the
left and right straight-boundary frames by
\[
 \mathcal B_{V,L}\ell_V=(-1,0),\quad \mathcal B_{V,L}(-p_V)=(0,1),
 \qquad
 \mathcal B_{V,R}r_V=(1,0),\quad \mathcal B_{V,R}(-p_V)=(0,1).
\]
The resulting matrices in the standard basis of $M^\vee$ are
\begin{equation}\label{eq:straight-boundary-frames}
\begin{array}{c|c|c}
 V&\mathcal B_{V,L}&\mathcal B_{V,R}\\ \hline
 A&\begin{pmatrix}-1&1\\0&1\end{pmatrix}
  &\begin{pmatrix}-1&1\\1&0\end{pmatrix}\\[5pt]
 B&\begin{pmatrix}1&1\\1&0\end{pmatrix}
  &\begin{pmatrix}1&1\\-1&-2\end{pmatrix}\\[5pt]
 C&\begin{pmatrix}-1&-3\\-1&-2\end{pmatrix}
  &\begin{pmatrix}-1&-3\\0&1\end{pmatrix}.
\end{array}
\end{equation}
The matrix $\check T_V=\mathcal B_{V,R}^{-1}\mathcal B_{V,L}$ is the monodromy of the tangent
lattice $M^\vee$ of $\Delta$ about $V$; it fixes $p_V$.  Along a
loop the dual lattice $M$ has the inverse transpose monodromy.  We orient
the loop $\gamma_V$ so that the monodromy of $M$ along it is
$T_V=\check T_V^T$, that is, oppositely to the loop along which the tangent
monodromy is $\check T_V$.  Direct multiplication yields precisely the
matrices \eqref{eq:disk-monodromies}.  Their invariant primitive characters
are, in standard coordinates of $M$, $\bar u_A=(1,-1)$, $\bar u_B=(1,1)$,
$\bar u_C=(1,3)$, and
\begin{equation}\label{eq:disk-rank-one}
\begin{aligned}
 T_A-\mathrm{Id}&=\binom{-1}{1}(1,1)=-\bar u_An_A^T,\\
 T_B-\mathrm{Id}&=\binom11(2,-2)=-2\,\bar u_Bn_B^T,\\
 T_C-\mathrm{Id}&=\binom13(3,-1)=-\bar u_Cn_C^T,
\end{aligned}
\end{equation}
with the primitive conormals $n_A=(1,1)$, $n_B=(-1,1)$, $n_C=(-3,1)$ in
$M^\vee$, in the coordinates $(s_1,s_2)$.  The conormals of
\eqref{eq:root-data} are
$N_{\widetilde A}=4n_A$, $N_{B_0}=4n_B$ and $N_{\widetilde C}=4n_C$: the
$F$-coordinates of $4n_A$, $4n_B$, $4n_C$ are $(1,-3)$, $(-1,-1)$, $(-3,1)$.
Thus the monodromies of the disk about $A$, $B$ and $C$ have widths $1$, $2$
and $1$.

In standard coordinates the exponent map and its inverse are
\[
 F=\begin{pmatrix}1/4&-1/4\\0&-1/2\end{pmatrix},
 \qquad
 F^{-1}=\begin{pmatrix}4&-2\\0&-2\end{pmatrix}.
\]
Conjugating by $F$ gives the disk monodromies in $F$-coordinates:
\begin{equation}\label{eq:unbranched-conjugates}
\begin{aligned}
 F^{-1}T_AF&=\begin{pmatrix}-1/2&9/2\\-1/2&5/2\end{pmatrix},\\
 F^{-1}T_BF&=\begin{pmatrix}2&1\\-1&0\end{pmatrix},\\
 F^{-1}T_CF&=\begin{pmatrix}-1/2&1/2\\-9/2&5/2\end{pmatrix}.
\end{aligned}
\end{equation}
The loops whose monodromy preserves $\widetilde M$ form a subgroup that
contains the parity subgroup, by \eqref{eq:parity-matrices} below, and does
not contain $\gamma_A$.  Since the parity subgroup has index two, it is
exactly the stabilizer of $\widetilde M$; in particular the monodromies about
$A$ and $C$ do not preserve $\widetilde M$.

For a loop $\gamma$ in the parity subgroup write $H_\gamma$ for the
monodromy of $\widetilde M$ along $\gamma$, with the convention
$H_{gh}=H_gH_h$, where a product $gh$ traverses $g$ first.  For the Schreier
transversal $\{1,\gamma_A\}$, the Reidemeister--Schreier basis of the parity
subgroup is
\[
 \gamma_B,\quad \gamma_C\gamma_A^{-1},\quad \gamma_A^2,\quad
 \gamma_A\gamma_B\gamma_A^{-1},\quad \gamma_A\gamma_C.
\]
The matrices of these loops, in $F$-coordinates, are
\begin{equation}\label{eq:parity-matrices}
 \begin{pmatrix}2&1\\-1&0\end{pmatrix},\quad
 \begin{pmatrix}-1&2\\-10&19\end{pmatrix},\quad
 \begin{pmatrix}-2&9\\-1&4\end{pmatrix},\quad
 \begin{pmatrix}-14&25\\-9&16\end{pmatrix},\quad
 \begin{pmatrix}-20&11\\-11&6\end{pmatrix}.
\end{equation}
All are integral.  The first is $H_{B,0}$; the second is the core matrix
$K=H_{\gamma_{\mathrm{core}}}$ of the loop
$\gamma_{\mathrm{core}}=\gamma_C\gamma_A^{-1}$; the third and fourth are
$H_A$ and $H_{B,1}$; the fifth is $H_{\gamma_A\gamma_C}=K^{-1}H_C$, of
trace $-14$.  Squaring the loop $\gamma_C$ gives
$H_C=H_{\gamma_C^2}=H_{\gamma_C\gamma_A^{-1}}H_{\gamma_A\gamma_C}$, the
final local matrix in \eqref{eq:lifted-monodromies}.  The primitive
factorizations \eqref{eq:primitive-factorizations} follow by subtracting the
identity.

The kinks of $\varphi$ across the bounded edges of the Legendre-dual
decomposition are, up to sign, the vertices of $\Delta$.  Equivalently, they
are the primitive facet conormals of $\Delta^\vee$, which in the coordinates
$(s_1,s_2)$ are $-p_A=(1,1)$, $-p_B=(1,-1)$ and $-p_C=(-3,1)$, each taking
the value $1$ on its facet.  Their $F$-coordinates are
\begin{equation}\label{eq:polarization-conormals}
 F^T(1,1)=(\tfrac14,-\tfrac34),\quad
 F^T(1,-1)=(\tfrac14,\tfrac14),\quad
 F^T(-3,1)=(-\tfrac34,\tfrac14).
\end{equation}
The kinks across the unbounded edges are the edge vectors of $\Delta$, which
are $(0,2)$, $(4,-2)$ and $(-4,0)$ in the coordinates $(s_1,s_2)$ and have
the integral $F$-coordinates $(0,-1)$, $(1,0)$ and $(-1,1)$.  Thus
$p^*\varphi$ is quarter-integral and $4p^*\varphi$ is the minimal integral
multiple.

\subsection{Core and boundary holonomy}

The core matrix satisfies
\[
 K-\mathrm{Id}=\begin{pmatrix}-2&2\\-10&18\end{pmatrix}.
\]
The lattice $M$ is generated by $g_1=F(4,0)$ and $g_2=F(-2,-2)$.  The
columns of $K-\mathrm{Id}$ are the $F$-coordinates of
$(K-\mathrm{Id})F(1,0)=2g_1+5g_2$ and $(K-\mathrm{Id})F(0,1)=-4g_1-9g_2$,
so $K$ acts trivially on $\widetilde M/M$.  It has trace
$18$ and determinant one, which proves the spectral statements in
Section~\ref{sec:jacobi-theta}.

The two boundary circles are the lifts of the boundary loop
$\gamma_A\gamma_B\gamma_C$, whose monodromy on $M$ is, in standard
coordinates, $T_AT_BT_C=\begin{pmatrix}1&0\\-8&1\end{pmatrix}$.  Their
monodromies on $\widetilde M$ are, in $F$-coordinates,
\begin{equation}\label{eq:boundary-matrices-app}
 U_0=H_{\gamma_A\gamma_B\gamma_C}=\begin{pmatrix}5&-4\\4&-3\end{pmatrix},
 \qquad
 U_1=H_{\gamma_A^2\gamma_B\gamma_C\gamma_A^{-1}}
 =\begin{pmatrix}33&-64\\16&-31\end{pmatrix},
\end{equation}
and $U_0=H_{B,1}K^{-1}H_C$, $U_1=H_AH_{B,0}K$.  They factor as
\[
 U_0-\mathrm{Id}=4\binom11(1,-1),
 \qquad
 U_1-\mathrm{Id}=16\binom21(1,-2).
\]
Using the invariant tangent--normal bases
\[
 P_0=\begin{pmatrix}1&0\\1&1\end{pmatrix},
 \qquad
 P_1=\begin{pmatrix}2&1\\1&1\end{pmatrix},
\]
one obtains
\[
 P_0^{-1}U_0P_0=\begin{pmatrix}1&-4\\0&1\end{pmatrix},
 \qquad
 P_1^{-1}U_1P_1=\begin{pmatrix}1&-16\\0&1\end{pmatrix}.
\]
This gives the two global widths used in Section~\ref{sec:compactification}.

On the side of $\Delta$, the invariant line of $\check T_V$ is the line
through $O$ and $p_V$, which contains $V$.  In coordinates centered at $O$
the affine transition at each singular point is therefore linear, and so is
the affine holonomy of every loop.  In the coordinates $(s_1,s_2)$, as in
Section~\ref{subsec:cofactor-supports}, the branch segment joins
$A=(-\frac12,-\frac12)$ to $C=(\frac32,-\frac12)$ and crosses the negative
half-chords at $(\frac23,-\frac12)$ and $(\frac13,-\frac12)$.  Its
interior lies in the complement of $c_A\cup c_B\cup c_C$, so its lifts lie
in the sheets $S_0,S_1$ of Section~\ref{subsec:kummer-charts}.

\subsection{Boundary development and exit slopes}

In the chart of the sheet $S_0$ one has
$\widetilde M^\vee=F^{-T}(\ZZ^2)=2L(\ZZ^2_{x,y})$.  The development below
is written in the plane of $\Delta$, in the invariant tangent--normal
coordinates $(\mathsf x,\mathsf y)$ of \eqref{eq:collar-holonomy},
normalized so that $\widetilde M^\vee$ is $2\ZZ^2$, equivalently so that the
lattice $L(\ZZ^2_{x,y})$ of cofactor directions is $\ZZ^2$.  The development
data of the two boundary circles are
\begin{equation}\label{eq:development-data}
\begin{array}{c|c|c|c}
 i&\ell_i&w_i&(\phi_i,h_i)\\ \hline
 0&4&4&(7/2,1)\\
 1&8&16&(15/2,1/2).
\end{array}
\end{equation}
Here $w_i$ is the width of the boundary monodromy, $\ell_i=w_ih_i$ is the
length of the boundary circle in these coordinates, and $(\phi_i,h_i)$ is
the common apex of the developed facets, in the coordinates
$(\mathsf x,\mathsf y)$.  Affine lengths measured in $\widetilde M^\vee$
are half the coordinate lengths; the widths and the slopes below do not
depend on this normalization.  In the coordinate $\mathsf r=-\mathsf y$, a
support of slope $\xi$ lies on the line
$\mathsf x=\phi_i-\xi(\mathsf r+h_i)$ of \eqref{eq:exit-line}.  On each line
$\{\mathsf r=\mathsf r_0\}$ the holonomy \eqref{eq:collar-holonomy} is the
translation by $w_i(\mathsf r_0+h_i)$, so its $k$th power maps the line of
slope $\xi$ to the line of slope $\xi-kw_i$.

By Proposition~\ref{prop:proper-supports} no support meets an arc carrying a
transition, so each support has direction $Lq$ in the chart of its sheet,
or $-Lq$ for a negative half-chord.  Write $q=c_1(2,1)+c_2(1,2)$ with
$c_1,c_2\geq0$, as in \eqref{eq:outgoing-cofactor}.  A support of
direction $-Lq$ leaves $\Delta$ through the facet $s_2=-1$, and one of
direction $Lq$ through the facet $s_1=-1$ if $c_1\geq c_2$ and through
$s_1+2s_2=1$ if $c_1\leq c_2$.  Developing the three facets along
$\partial_0\overline\Acal$ into the coordinates above gives the slopes
\[
 \xi_-=\frac{2c_1+c_2}{3(c_1+c_2)},
 \qquad
 \xi_+=
 \begin{cases}
 \dfrac{3(7c_1+3c_2)}{4(2c_1+c_2)},&c_1\geq c_2,\\[10pt]
 \dfrac{11c_1+19c_2}{4(c_1+2c_2)},&c_1\leq c_2,
 \end{cases}
\]
for the directions $-Lq$ and $Lq$; along $\partial_1\overline\Acal$ the
slopes are $4\xi_\pm+1$.  We call the slopes of supports of direction
$-Lq$ and $Lq$ the \emph{negative} and \emph{positive slopes}, after the
signs of these directions; as numbers, all of them are positive.  Each of
the functions above is decreasing in $c_2/c_1\in[0,\infty]$, and the two
expressions for $\xi_+$ agree at $c_1=c_2$.  The negative and positive
slopes therefore lie in the closed intervals bounded by their values at
$c_2=0$ and at $c_1=0$:
\begin{equation}\label{eq:slope-envelopes}
\begin{array}{c|c|c}
 i&\text{negative slopes}&\text{positive slopes}\\ \hline
 0&[1/3,2/3]&[19/8,21/8]\\
 1&[7/3,11/3]&[21/2,23/2].
\end{array}
\end{equation}
For $i=0$, every difference between a positive and a negative slope lies in
$(0,4)$, and for $i=1$ it lies in $(0,16)$; each interval has width smaller
than $w_i$.  Thus the slopes of two supports never differ by a nonzero
multiple of $w_i$, and no nonzero power of the holonomy
\eqref{eq:collar-holonomy} maps a support onto a support.  These are the
inequalities used in the proof of
Proposition~\ref{prop:exit-separation}.

\subsection{Finite root algebra}

For an effective Cartier divisor $D_f=(f)$ on $X=\Spec A$, we call
$\Spec A[g]/(g^d-f)$ the \emph{Kummer cover} of $(X,f)$, with $\mu_d$ acting
by $g\mapsto\zeta g$ and trivially on $A$.  Our root-stack convention is
\begin{equation}\label{eq:root-stack-chart-app}
 \sqrt[d]{(X,D_f)}
 =\left[\Spec A[g]/(g^d-f)\big/\mu_d\right].
\end{equation}
Over $X\setminus D_f$ the Kummer cover is the $\mu_d$-torsor
$\Spec A[f^{-1}][g]/(g^d-f)$, finite étale because $d$ is invertible in
$\Bbbk$, and the root stack is $X\setminus D_f$ itself: its functions there
are the invariants $A[f^{-1}]$, and $g$, of weight one, is not among them.
Along $D_f$ the quotient retains the root automorphism group, and $g=0$
defines the root divisor.  If $A$ is a character algebra and $N$ an integral
conormal vanishing on the exponents of $f$, the slab map
$z^m\mapsto g^{\langle N,m\rangle}z^m$ is an automorphism of the Kummer ring
$A[f^{-1}][g]/(g^d-f)$ fixing $f$ and $g$.  It commutes with the action of
$\mu_d$ on $z^m$ exactly when $d$ divides $\langle N,m\rangle$.  It
therefore acts on the Kummer cover over $X\setminus D_f$, and it descends to
the root stack only if $d$ divides $\langle N,m\rangle$ for every character
$m$.  For divisors $D_e=(f_e)$ with degrees $d_e$, $e=1,\dots,s$, the
fiber product over $X$ of the root stacks $\sqrt[d_e]{(X,D_e)}$ is the
quotient of the Kummer cover $\Spec A[g_e]_e/(g_e^{d_e}-f_e)_e$ by
$\prod_e\mu_{d_e}$, and the stabilizer of a point of this cover is the
product of the factors $\mu_{d_e}$ whose coordinate $g_e$ vanishes at the
point, that is, of the factors with $x\in D_e$ for the image $x$ of the
point in $X$.  Section~\ref{sec:compactification}
applies this to the finite root stacks $\Xfrak_T$ of
Section~\ref{sec:affine-kummer}.

Let $\mathfrak I$ be the filtered poset of finite rooted subtrees of the path
tree $\Tcal$.  For $T\subset T'$ in $\mathfrak I$, the ring
$R_{T'}^{\mathrm{root}}$ of \eqref{eq:finite-kummer-ring} is free over
$\mathcal O(U_{T'})$ on the monomials in the roots with exponents smaller
than their degrees, and these contain the monomials of $T$; hence
$R_T^{\mathrm{root}}\to R_{T'}^{\mathrm{root}}$ is injective.  The colimit
over $\mathfrak I$ is the pathwise Kummer ring
$R_{\mathrm{path}}^{\mathrm{root}}$.  A finitely presented module, algebra,
morphism, or finite collection of identities among composites over
$R_{\mathrm{path}}^{\mathrm{root}}\otimes_\Bbbk R/\mathfrak m^N$ uses finitely many
generators, relations, and matrix entries.  All therefore occur over
$R_T^{\mathrm{root}}\otimes_\Bbbk R/\mathfrak m^N$ for one $T\in\mathfrak I$, after
replacing the chosen stages by their rooted union.  This proves the
finite-presentation assertion in Theorem~\ref{thm:forced-category}.
